\section{Uniform minimax rates across noise scales}

The public certificate in \cref{obj:thm:finite-certificate} bounds both estimation error and honest interval length. We now characterize the optimal scale of these two criteria across sample sizes and calibrated noise levels. The characterization uses a scalar resolution equation that balances endpoint approximation against the cost of Gaussian inversion. Its definition also identifies the observable procedures that attain the resulting rate.

\begin{definitionv}{def:uniform-rate}[Uniform rate \(r_\beta(n,\sigma)\)]
The uniform resolution rate and selected observable procedures are
\[
r_\beta(n,\sigma)=h_\star(n,\sigma)^\beta,
\qquad
\widetilde\theta_{\beta,n,\sigma}=\widehat\theta_{L_\star},
\qquad
\widetilde I_{\beta,n,\sigma}=I_{L_\star},
\]
where \(L_\star\) is selected by the finite public minimization. With the local abbreviations
\[
\nu=2\beta+1,\qquad h_0=n^{-1/\nu},
\]
define the noise cost, for \(0<h\le1\), by
\[
\Psi(h,\sigma)=
\begin{cases}
0,&\sigma\le h,\\
(\sigma/h)^{2/3}-1,&\sigma^4\le h<\sigma,\\
h^{-1/2}\{1+\log(\sigma^2/\sqrt h)\}-1,&h<\sigma^4.
\end{cases}
\]
At \(\sigma=0\), the first branch applies throughout \(0<h\le1\). The resolution \(h_\star(n,\sigma)\) is specified as the root in \([h_0,1]\) of
\[
\log n+\nu\log h_\star(n,\sigma)
=\Psi(h_\star(n,\sigma),\sigma).
\]
The left side is continuous and increasing in the resolution, and the right side is continuous and nonincreasing. Root existence and uniqueness, risk bounds, and interval honesty are subsequent assertions. All tuning parameters belong to the procedures; the law class retains its defining restrictions.
\end{definitionv}

The noise cost \leanref{sym:Noisecost}{\ensuremath{\Psi(h,\sigma)}} summarizes Gaussian inversion cost for the rate comparison. The selected resolution \leanref{sym:Rateresolution}{\ensuremath{h_\star(n,\sigma)}} translates this cost into an endpoint approximation scale, and the rate \leanref{sym:Frontier}{\ensuremath{r_\beta(n,\sigma)}} is its Hölder power. The selected estimator \leanref{sym:Attainer}{\ensuremath{\widetilde\theta_{\beta,n,\sigma}}} and selected interval \leanref{sym:Honestattainer}{\ensuremath{\widetilde I_{\beta,n,\sigma}}} are the procedures chosen by \cref{obj:def:public-selection}. Thus the scalar equation describes the precision achieved by the finite public selection rule.

The two noisy branches come from distinct parts of the same endpoint mechanism. In the intermediate range, the noise cost grows as \((\sigma/h)^{2/3}-1\), linking finer endpoint resolution to a larger Gaussian inversion cost. Below the compact-support interface \(h=\sigma^4\), the cost becomes \(h^{-1/2}\{1+\log(\sigma^2/\sqrt h)\}-1\). These expressions agree at the interface. The balance equation therefore connects endpoint approximation to inversion cost continuously across the two noisy regimes, using the same resolution to compare both decision criteria.

On the benchmark class of \cref{obj:def:model}, the following result characterizes minimax absolute risk and minimax honest expected length simultaneously. Its comparison constants are uniform in both sample size and noise scale.

\begin{theoremv}{thm:uniform-frontier}[Uniform estimation and interval frontier]
Fix \(\beta\in(0,1]\). There exist constants \(c,C>0\), depending only on \(\beta\), such that the following conclusions hold simultaneously for every \(n\ge2\) and every \(\sigma\in[0,1]\).

Write \(\nu=2\beta+1\) and \(h_0=n^{-1/\nu}\). The equation defining the resolution \(h_\star(n,\sigma)\) has a unique root. With \(r_\beta(n,\sigma)=h_\star(n,\sigma)^\beta\), the selected observable estimator \(\widetilde\theta_{\beta,n,\sigma}\), the selected observable interval \(\widetilde I_{\beta,n,\sigma}\), and their public radius \(E_\beta(n,\sigma)\) satisfy
\[
c r_\beta(n,\sigma)
\le R_\beta(n,\sigma)
\le \sup_{P\in\mathcal M_\beta(\sigma)}
\mathbb E_{P,n}
\left|\widetilde\theta_{\beta,n,\sigma}-\theta(P)\right|
\le E_\beta(n,\sigma)
\le C r_\beta(n,\sigma)
\]
and
\[
c r_\beta(n,\sigma)
\le \mathcal L_\beta(n,\sigma)
\le \sup_{P\in\mathcal M_\beta(\sigma)}
\mathbb E_{P,n}
\operatorname{length}(\widetilde I_{\beta,n,\sigma})
\le 2E_\beta(n,\sigma)
\le C r_\beta(n,\sigma),
\qquad
\widetilde I_{\beta,n,\sigma}\in\mathcal I_{\beta,n,\sigma}.
\]
The minimax criteria range over all procedures in their stated decision classes, including procedures with arbitrary independent randomization and arbitrary degree.

The following branches exhaust every sample size. Throughout, \(\asymp\) denotes comparison by positive constants depending only on \(\beta\), except where dependence on an additional fixed exponent is stated.
\begin{itemize}
\item \textbf{(Direct branch.)}
If \(\sigma\le h_0\), then \(h_\star=h_0\), including at \(\sigma=0\).

\item \textbf{(Intermediate branch.)}
If \(\sigma>h_0\), then \(h_0<h_\star<\sigma\). Within this noisy region, the intermediate branch occurs exactly when
\[
\log(n\sigma^{4\nu})\le \sigma^{-2}-1.
\]
There is a unique \(z\in(1,\sigma^{-2}]\) satisfying
\[
z+\frac{3\nu}{2}\log z=1+\log(n\sigma^\nu),
\]
and
\[
h_\star=\sigma z^{-3/2}
\asymp
\frac{\sigma}{[1+\log(n\sigma^\nu)]^{3/2}}.
\]

\item \textbf{(Compact branch.)}
If \(\sigma>h_0\) and
\[
\log(n\sigma^{4\nu})>\sigma^{-2}-1,
\]
there is a unique \(\tau>\sigma^{-2}\) satisfying
\[
2\nu\log\tau+
\tau\{1+\log(\sigma^2\tau)\}
=\log n+1.
\]
On this branch,
\[
h_\star=\tau^{-2},
\qquad
\tau\asymp
\frac{\log(en)}{\log(e+\sigma^2\log(en))},
\qquad
r_\beta(n,\sigma)\asymp
\left\{
\frac{\log(e+\sigma^2\log(en))}{\log(en)}
\right\}^{2\beta},
\]
uniformly over the branch.
\end{itemize}
At the intermediate--compact equality, \(h_\star=\sigma^4\), and both scalar formulas give this value. At the direct--noisy equality, \(h_\star=\sigma=h_0\), and the adjacent formulas agree.

For every fixed \(a\) with \(0<a<1/(2\beta+1)\),
\[
r_\beta(n,n^{-a})
\asymp n^{-a\beta}(\log n)^{-3\beta/2}
\]
as \(n\) tends to infinity, with comparison constants that may additionally depend on \(a\). For every fixed \(a\ge1/(2\beta+1)\), the sequence \(\sigma=n^{-a}\) lies in the exact direct branch for every \(n\ge2\). For every fixed \(\sigma\in(0,1]\),
\[
r_\beta(n,\sigma)
\asymp
\left\{\frac{\log\log(en)}{\log(en)}\right\}^{2\beta}
\]
as \(n\) tends to infinity, with comparison constants allowed to depend on the fixed noise level. The root and finite-sample branch formulas characterize every shrinking-noise sequence.

In particular, for \(\beta=1\) and \(\sigma=n^{-1/6}\), the intermediate branch holds for every \(n\ge2\), and
\[
r_1(n,n^{-1/6})
=h_\star(n,n^{-1/6})
\asymp
\frac{n^{-1/6}}{[1+\tfrac12\log n]^{3/2}}.
\]
\end{theoremv}

\Cref{obj:thm:uniform-frontier} contains the exact branch conditions, both interfaces, and the relevant shrinking- and fixed-noise specializations. Its scalar balance connects endpoint approximation with Gaussian inversion cost, while the matched upper and lower bounds show that this one resolution governs both absolute risk and honest expected length. In the direct regime, the inversion cost vanishes at the direct resolution. In the intermediate regime, resolving the endpoint more finely than the noise scale incurs the fractional-power cost in \cref{obj:def:uniform-rate}. The compact-support regime uses the finer-resolution cost below \(\sigma^4\), yielding the logarithmic comparison in \cref{obj:thm:uniform-frontier}.

For reference, \cref{tab:noise-regimes} collects the exact branch conditions and resolutions. The auxiliary scalars \(z\) and \(\tau\) are the unique solutions specified in \cref{obj:thm:uniform-frontier}.

\begin{table}[htbp]
\centering
\caption{Exact resolution regimes for \(n\ge2\) and \(\sigma\in[0,1]\), with \(\nu=2\beta+1\) and \(h_0=n^{-1/\nu}\).}
\label{tab:noise-regimes}
\begin{tabular}{@{}p{0.17\linewidth}p{0.49\linewidth}p{0.25\linewidth}@{}}
\hline
Regime & Exact condition & Exact resolution \\
\hline
Direct &
\(\sigma\le h_0\) &
\(h_\star=h_0\) \\[0.6em]
Intermediate &
\(\sigma>h_0\) and
\(\log(n\sigma^{4\nu})\le\sigma^{-2}-1\) &
\(h_\star=\sigma z^{-3/2}\) \\[0.6em]
Compact support &
\(\sigma>h_0\) and
\(\log(n\sigma^{4\nu})>\sigma^{-2}-1\) &
\(h_\star=\tau^{-2}\) \\
\hline
\end{tabular}

\smallskip
\parbox{0.95\linewidth}{\small
At the direct--noisy interface, \(\sigma=h_0\) and \(h_\star=\sigma=h_0\).
At the intermediate--compact interface,
\(\log(n\sigma^{4\nu})=\sigma^{-2}-1\) within the noisy region, and
\(h_\star=\sigma^4\); both noisy scalar formulas agree there.
The common rate is \(h_\star^\beta\), and its comparison with both minimax criteria uses constants depending only on \(\beta\), uniformly in \(n\) and \(\sigma\).}
\end{table}

The uniform characterization also distinguishes asymptotic paths. For polynomially shrinking noise \(\sigma=n^{-a}\), a fixed exponent \(0<a<1/(2\beta+1)\) gives the order \(n^{-a\beta}(\log n)^{-3\beta/2}\), whereas \(a\ge1/(2\beta+1)\) places every sample size in the exact direct branch. With fixed positive noise, the rate has order \(\{\log\log(en)/\log(en)\}^{2\beta}\). The Lipschitz example \(\beta=1\), \(\sigma=n^{-1/6}\) sharpens this comparison: the intermediate condition holds for every \(n\ge2\), so its displayed formula describes the entire sequence. These are comparisons of rates; the exact resolution is determined by the scalar equation, and the minimax inequalities retain their comparison constants. Constants for the uniform branch comparisons depend only on \(\beta\), while those for a specified asymptotic path may also depend on its fixed exponent or fixed noise level as stated in \cref{obj:thm:uniform-frontier}.

At zero noise, endpoint approximation and sampling variability balance at \(h_0=n^{-1/(2\beta+1)}\), giving the Hölder power \(n^{-\beta/(2\beta+1)}\). The direct-experiment specialization in \cref{obj:thm:direct-reduction}, presented in the verification appendix, establishes this order for both criteria and its attainment by the selected estimator and honest interval. Together with the exact direct branch in \cref{obj:thm:uniform-frontier}, it shows how sufficiently precise score measurement recovers the direct-experiment rate throughout the region \(\sigma\le h_0\).
